% ID: SEL-D0011-BD-E07
% ESTADO: cotejada. Segunda pasada visual completa y borrador compilado revisados.
% Fuente: Selectividad/matematicas_aplicadas/EXTRA_SUPLENTE_Modelo-1-Enunciados (Mat-Ap-2021).pdf
% SHA-256: 4a6c7d34ff35c3e7678c52970b4d4558cee2c35224b46c66035cb0bc8d5a7ec8
% Página: 2; bloque D; ejercicio original 7.

\begin{enumerate}[label=\alph*)]
  \item \textbf{(1 punto)} En una población constituida por los números
  naturales del 1 al 9, \textquestiondown{}cuántas muestras de tamaño 2 se
  pueden formar por muestreo aleatorio simple? Si se elige al azar una de esas
  muestras, \textquestiondown{}cuál es la probabilidad de que el valor medio
  de los dos números de esa muestra sea 5?
  \item \textbf{(1.5 puntos)} Para estimar la proporción de andaluces contagiados
  por una enfermedad infecciosa en un momento determinado, se ha tomado una
  muestra de 10\,000 personas, resultando que 500 de ellas estaban infectadas.
  \begin{enumerate}[label=\arabic*.]
    \item Con ese dato, establezca un intervalo, al 97\% de confianza, para
    la proporción real de infectados en la población andaluza.
    \item A la vista del intervalo obtenido, razone si se podría aceptar que
    el 6\% de la población andaluza estaba infectada.
    \item Se toma una nueva muestra de mayor tamaño y resulta que hay la
    misma proporción de positivos en la nueva muestra. Con estos nuevos datos,
    razone si el nuevo intervalo al 97\% de confianza contiene al intervalo
    anterior o está contenido en él.
  \end{enumerate}
\end{enumerate}
